\documentclass[12pt]{article}
\usepackage{graphicx}
\usepackage{booktabs}
\usepackage{geometry}
\usepackage{hyperref}
\usepackage{amsmath}
\geometry{margin=1in}

\title{
Optical Information Persistence in Historical Oil Portraits:\\
A BIMOD-Based Comparative Polarity Study of Mona Lisa and King George IV
}

\author{
DuckJin Yoon (Solvaram)\\
Independent Researcher\\
OSF Project DOI: 10.17605/OSF.IO/UBQKE
}

\date{February 2026}

\begin{document}
\maketitle

\begin{abstract}
This study reports reproducible polarity-specific responses obtained using the Bodily Interferometric Magneto-Optical Discrimination (BIMOD) protocol when applied to historical oil portraits. High-resolution reproductions of Leonardo da Vinci’s Mona Lisa and Thomas Lawrence’s portrait of King George IV were examined. Despite differences in artistic style, historical period, and subject identity, both paintings elicited consistent polarity-dependent resonance responses under standardized experimental conditions. The presence of reproducible polarity reactions suggests preservation of structured optical information within these historical images. No mechanistic explanation is claimed.
\end{abstract}

\section{Introduction}

The BIMOD protocol was introduced as a human-response-based polarity discrimination method involving biological magnetic nanoparticle resonance \cite{Yoon2026BIMOD}. While initially developed for biological polarity detection, its application to visual images has revealed reproducible responses associated with specific structured stimuli.

Historical oil portraiture, particularly in works associated with optical realism and tonal depth layering, presents an opportunity to investigate whether structured optical information may persist within painted surfaces.

This study does not attempt reinterpretation of art history. It documents reproducible polarity-specific responses observed under controlled procedural conditions.

\section{Materials and Methods}

\subsection{Paintings Examined}

Two classical portraits were selected:

\begin{itemize}
\item Mona Lisa — Leonardo da Vinci (c.1503–1506)
\item King George IV — Thomas Lawrence
\end{itemize}

High-resolution digital reproductions were used under standardized lighting.

\subsection{BIMOD Protocol}

A cylindrical bar magnet (~10 cm length) was held vertically with the lower N-pole end grasped in the left hand. The right arm remained fully relaxed.

A positive response was defined as spontaneous arm elevation occurring only when image polarity was opposite to the handheld magnet polarity.

Measurements were repeated multiple times to confirm reproducibility.

Control images lacking structured tonal gradation were also tested.

\section{Results}

\begin{table}[h]
\centering
\begin{tabular}{lcc}
\toprule
Parameter & Mona Lisa & King George IV \\
\midrule
Nasion Code & 1S & 2N \\
BIMOD Response & Arm elevation & Arm elevation \\
Reproducibility & Confirmed & Confirmed \\
\bottomrule
\end{tabular}
\caption{Comparative Polarity Response}
\end{table}

Both portraits produced polarity-specific resonance responses consistent with predefined BIMOD criteria.

Control images did not produce reproducible responses.

\section{Discussion}

The two portraits differ in historical period, painter, stylistic execution, and subject identity. Despite these differences, consistent polarity responses were observed.

The presence of reproducible polarity reactions suggests preservation of structured optical information within the painted image.

This study does not claim a mechanistic explanation. It presents observational evidence intended to motivate interdisciplinary research into optical encoding, perceptual physics, and structured image information retention.

\section{Instrumentation Development Plan}

Human-response-based polarity detection must transition toward objective measurement systems. Ongoing efforts focus on identifying micro-voltage (µV) scale biosignal sensors capable of detecting transient resonance amplification during BIMOD exposure.

Future studies will integrate instrumental measurement to quantify polarity-specific responses independently of subjective arm-elevation criteria.

\section{Conclusion}

Reproducible polarity-specific responses were observed in two historical oil portraits using the BIMOD protocol. The findings suggest structured optical information persistence and warrant further controlled investigation.

\bibliographystyle{plain}
\bibliography{BIMOD_100_references}

\end{document}
